\documentclass{beamer}
\usetheme{metropolis}
\usepackage{graphicx}
\usepackage{booktabs}
\usepackage{amsmath,amssymb}
\usepackage{CJKutf8}
\graphicspath{ {D:/GitHub/PaperHub/backend/workspace/chat_session/172/slides/figures/} }
\title{Action-aware Dynamic Pruning for Efficient Vision-Language-Action Manipulation}
\subtitle{針對視覺-語言-動作模型之動作感知動態剪枝技術}
\author{Pei, Chen, Xu, et al.}
\date{arXiv:2509.22093v1 (2025)}
\begin{document}
\begin{CJK*}{UTF8}{bsmi}
\begin{frame}[plain]
\titlepage
\end{frame}
\begin{frame}
\frametitle{Executive Summary}
\begin{itemize}
    \item Proposes Action-aware Dynamic Pruning (ADP) for Vision-Language-Action (VLA) models.
    \item Toggles visual token pruning based on end-effector trajectories.
    \item Prunes during coarse movements and retains context for fine manipulations.
    \item Achieves 1.35x speedup and 25.8\% higher task success rate.
\end{itemize}
\end{frame}
\begin{frame}
\frametitle{Computational Bottlenecks in VLA Models}
\begin{itemize}
    \item High-resolution visual inputs generate massive token sequences.
    \item Excessive FLOPs cause severe inference latency.
    \item Closed-loop robotic control requires high-frequency responses.
    \item Static pruning models fail strict latency and accuracy constraints.
\end{itemize}
\end{frame}
\begin{frame}
\frametitle{Phased Characteristics and Visual Redundancy}
\begin{itemize}
    \item Manipulation tasks divide into coarse and fine-grained phases.
    \item Coarse movements exhibit high visual redundancy and fault tolerance.
    \item Fine-grained operations require precise visual details for success.
    \item Aggressive static pruning destroys context during fine manipulation.
\end{itemize}
\begin{center}
    \includegraphics[width=0.85\textwidth,height=0.35\textheight,keepaspectratio]{p0-fig-001}
\end{center}
\end{frame}
\begin{frame}
\frametitle{Overview of Action-Aware Dynamic Pruning (ADP)}
\begin{itemize}
    \item Features Action-aware Gating and Anticipatory Pruning mechanisms.
    \item Dynamically triggers pruning based on recent end-effector trajectories.
    \item Filters task-relevant visual features using text-to-vision cross-attention.
\end{itemize}
\begin{center}
    \includegraphics[width=0.85\textwidth,height=0.4\textheight,keepaspectratio]{p0-fig-002}
\end{center}
\end{frame}
\begin{frame}
\frametitle{Attention-Based Visual Token Selection}
\begin{itemize}
    \item Assesses visual relevance via the VLA cross-attention mechanism.
    \item Scores identify critical regions for executing current instructions.
    \item Ranks and retains top $K$ tokens to discard background redundancy.
\end{itemize}
\[
\Phi^{(l)}(v) = \frac{1}{N^h \cdot L_{\text{txt}}} \sum_{h=1}^{N^h} \sum_{t=1}^{L_{\text{txt}}} \mathbf{A}^{(l)}_{h,t,v}
\]
\begin{equation*}
\mathbf{X}_{\text{keep}} = \operatorname{Top-K}\big(\Phi^{(l)}, \lfloor \rho \cdot L^{\text{vis}} \rfloor\big)
\end{equation*}
\end{frame}
\begin{frame}
\frametitle{Action-Aware Gating Mechanism}
\begin{itemize}
    \item Computes cumulative trajectory distance $\delta_i$ in observation windows.
    \item High distance denotes fast movement, triggering aggressive visual pruning.
    \item Dual thresholds establish hysteresis to prevent unstable state toggling.
\end{itemize}
\begin{equation*}
\delta_i=\sum_{t=b_i}^{e_i-1}\lVert \mathbf{p}_{t+1}-\mathbf{p}_t\rVert_2
\end{equation*}
\begin{equation*}
s_{i+1} = \begin{cases} 1, & \delta_i \ge U^{(i)}, \\ 0, & \delta_i \le V^{(i)}, \\ s_i, & V^{(i)} < \delta_i < U^{(i)}. \end{cases}
\end{equation*}
\end{frame}
\begin{frame}
\frametitle{Experimental Results: Simulation Performance}
\begin{itemize}
    \item Evaluated ADP against baselines on the LIBERO simulation suite.
    \item Overcomes real-time bottlenecks, achieving a 1.35x overall speedup.
    \item Filters visual noise effectively during early movement phases.
    \item Achieved up to 25.8\% higher success rate over the OpenVLA baseline.
\end{itemize}
\end{frame}
\begin{frame}
\frametitle{Experimental Results: Real-World Deployment}
\begin{itemize}
    \item Deployed ADP on Jaco2 hardware for diverse grasping tasks.
    \item Highly robust against complex lighting and dynamic disturbances.
    \item Low-latency inference crucially improves overall robot control smoothness.
\end{itemize}
\begin{center}
    \includegraphics[width=0.85\textwidth,height=0.45\textheight,keepaspectratio]{p0-fig-004}
\end{center}
\end{frame}
\begin{frame}
\frametitle{Conclusion and Future Work}
\begin{itemize}
    \item ADP effectively balances computational efficiency and manipulation accuracy.
    \item Successfully integrates visual attention mechanisms with robot kinematics.
    \item Solves real-time constraints in VLA models without hardware modifications.
    \item Future work will explore multimodal dynamic fusion using tactile feedback.
\end{itemize}
\end{frame}
\end{CJK*}
\end{document}